\hypertarget{huxe9ritage}{%
\section{Héritage}\label{huxe9ritage}}

\hypertarget{objets-carte}{%
\subsection{Objets Carte}\label{objets-carte}}

Il y a cinquante-deux cartes dans un jeu, dont chacune appartient à
l\textquotesingle une des quatre couleurs et l\textquotesingle une des
treize valeurs. Pour représenter une carte à jouer, les attributs
évidents sont la valeur (\emph{rank}) et la couleur (\emph{suit}). Une
alternative pour les représenter consiste à utiliser des entiers pour
encoder les valeurs et les couleurs, ce qui permet de comparer
facilement les cartes.

La définition de la classe \texttt{Card} ressemble à ceci :

\begin{Shaded}
\begin{Highlighting}[]
\KeywordTok{class}\NormalTok{ Card(}\BuiltInTok{object}\NormalTok{):}
    \CommentTok{"""Représente une carte à jouer standard."""}

    \KeywordTok{def} \FunctionTok{\_\_init\_\_}\NormalTok{(}\VariableTok{self}\NormalTok{, suit}\OperatorTok{=}\DecValTok{0}\NormalTok{, rank}\OperatorTok{=}\DecValTok{2}\NormalTok{):}
        \VariableTok{self}\NormalTok{.suit }\OperatorTok{=}\NormalTok{ suit}
        \VariableTok{self}\NormalTok{.rank }\OperatorTok{=}\NormalTok{ rank}
\end{Highlighting}
\end{Shaded}

\hypertarget{attributs-de-classe}{%
\subsection{Attributs de classe}\label{attributs-de-classe}}

Pour afficher les objets \texttt{Card} d\textquotesingle une manière
lisible, nous pouvons utiliser des attributs de classe définis à
l\textquotesingle intérieur de la classe mais en dehors de toute méthode
:

\begin{Shaded}
\begin{Highlighting}[]
\CommentTok{\# à l\textquotesingle{}intérieur de la classe Card :}

\NormalTok{suit\_names }\OperatorTok{=}\NormalTok{ [}\StringTok{"Clubs"}\NormalTok{, }\StringTok{"Diamonds"}\NormalTok{, }\StringTok{"Hearts"}\NormalTok{, }\StringTok{"Spades"}\NormalTok{]}
\NormalTok{rank\_names }\OperatorTok{=}\NormalTok{ [}
    \VariableTok{None}\NormalTok{,}
    \StringTok{"Ace"}\NormalTok{,}
    \StringTok{"2"}\NormalTok{,}
    \StringTok{"3"}\NormalTok{,}
    \StringTok{"4"}\NormalTok{,}
    \StringTok{"5"}\NormalTok{,}
    \StringTok{"6"}\NormalTok{,}
    \StringTok{"7"}\NormalTok{,}
    \StringTok{"8"}\NormalTok{,}
    \StringTok{"9"}\NormalTok{,}
    \StringTok{"10"}\NormalTok{,}
    \StringTok{"Jack"}\NormalTok{,}
    \StringTok{"Queen"}\NormalTok{,}
    \StringTok{"King"}\NormalTok{,}
\NormalTok{]}


\KeywordTok{def} \FunctionTok{\_\_str\_\_}\NormalTok{(}\VariableTok{self}\NormalTok{):}
    \ControlFlowTok{return} \StringTok{"}\SpecialCharTok{\%s}\StringTok{ of }\SpecialCharTok{\%s}\StringTok{"} \OperatorTok{\%}\NormalTok{ (Card.rank\_names[}\VariableTok{self}\NormalTok{.rank], Card.suit\_names[}\VariableTok{self}\NormalTok{.suit])}
\end{Highlighting}
\end{Shaded}

Les variables comme \texttt{suit\_names} et \texttt{rank\_names} sont
appelées des \textbf{attributs de classe} car elles sont associées à
l\textquotesingle objet classe \texttt{Card}, tandis que \texttt{suit}
et \texttt{rank} sont des attributs d\textquotesingle instance.

\hypertarget{comparaison-de-cartes}{%
\subsection{Comparaison de cartes}\label{comparaison-de-cartes}}

Pour les types intégrés, les opérateurs relationnels comparent les
valeurs. Pour les types définis par l\textquotesingle utilisateur, nous
pouvons redéfinir ce comportement en fournissant une méthode nommée
\texttt{\_\_cmp\_\_} qui prend deux paramètres, \texttt{self} et
\texttt{other}.

En Python 3, la méthode \texttt{\_\_cmp\_\_} n\textquotesingle est pas
prise en charge et il convient plutôt de fournir \texttt{\_\_lt\_\_}.

\hypertarget{jeux-de-cartes}{%
\subsection{Jeux de cartes}\label{jeux-de-cartes}}

Puisqu\textquotesingle un jeu de cartes (\emph{deck}) est composé de
cartes, il est naturel qu\textquotesingle il contienne une liste de
cartes comme attribut. La méthode d\textquotesingle initialisation crée
l\textquotesingle attribut \texttt{cards} et génère les cinquante-deux
cartes standard :

\begin{Shaded}
\begin{Highlighting}[]
\KeywordTok{class}\NormalTok{ Deck(}\BuiltInTok{object}\NormalTok{):}

    \KeywordTok{def} \FunctionTok{\_\_init\_\_}\NormalTok{(}\VariableTok{self}\NormalTok{):}
        \VariableTok{self}\NormalTok{.cards }\OperatorTok{=}\NormalTok{ []}
        \ControlFlowTok{for}\NormalTok{ suit }\KeywordTok{in} \BuiltInTok{range}\NormalTok{(}\DecValTok{4}\NormalTok{):}
            \ControlFlowTok{for}\NormalTok{ rank }\KeywordTok{in} \BuiltInTok{range}\NormalTok{(}\DecValTok{1}\NormalTok{, }\DecValTok{14}\NormalTok{):}
\NormalTok{                card }\OperatorTok{=}\NormalTok{ Card(suit, rank)}
                \VariableTok{self}\NormalTok{.cards.append(card)}
\end{Highlighting}
\end{Shaded}

\hypertarget{ajout-suppression-muxe9lange-et-tri}{%
\subsection{Ajout, suppression, mélange et
tri}\label{ajout-suppression-muxe9lange-et-tri}}

Pour distribuer des cartes, nous pouvons utiliser la méthode de liste
\texttt{pop} pour retirer une carte du jeu :

\begin{Shaded}
\begin{Highlighting}[]
\CommentTok{\# à l\textquotesingle{}intérieur de la classe Deck :}


\KeywordTok{def}\NormalTok{ pop\_card(}\VariableTok{self}\NormalTok{):}
    \ControlFlowTok{return} \VariableTok{self}\NormalTok{.cards.pop()}
\end{Highlighting}
\end{Shaded}

Pour ajouter une carte, nous pouvons utiliser la méthode \texttt{append}
:

\begin{Shaded}
\begin{Highlighting}[]
\CommentTok{\# à l\textquotesingle{}intérieur de la classe Deck :}


\KeywordTok{def}\NormalTok{ add\_card(}\VariableTok{self}\NormalTok{, card):}
    \VariableTok{self}\NormalTok{.cards.append(card)}
\end{Highlighting}
\end{Shaded}

Une méthode de ce type qui utilise une autre fonction sans effectuer de
travail complexe est parfois appelée un \textbf{placage}
(\emph{veneer}). Nous pouvons également écrire une méthode de mélange
(\emph{shuffle}) en utilisant la fonction \texttt{shuffle} du module
\texttt{random} :

\begin{Shaded}
\begin{Highlighting}[]
\CommentTok{\# à l\textquotesingle{}intérieur de la classe Deck :}


\KeywordTok{def}\NormalTok{ shuffle(}\VariableTok{self}\NormalTok{):}
\NormalTok{    random.shuffle(}\VariableTok{self}\NormalTok{.cards)}
\end{Highlighting}
\end{Shaded}

\hypertarget{huxe9ritage-1}{%
\subsection{Héritage}\label{huxe9ritage-1}}

La fonctionnalité de langage la plus souvent associée à la programmation
orientée objet est l\textquotesingle{}\textbf{héritage}, qui est la
capacité de définir une nouvelle classe constituant une version modifiée
d\textquotesingle une classe existante. La classe existante est appelée
la \textbf{classe parente} et la nouvelle classe est appelée la
\textbf{classe enfant} ou sous-classe.

Par exemple, pour représenter une « main » de cartes détenue par un
joueur, nous pouvons créer une classe \texttt{Hand} qui hérite de
\texttt{Deck} :

\begin{Shaded}
\begin{Highlighting}[]
\KeywordTok{class}\NormalTok{ Hand(Deck):}
    \CommentTok{"""Représente une main de cartes à jouer."""}
\end{Highlighting}
\end{Shaded}

Cette définition indique que \texttt{Hand} hérite de \texttt{Deck}, ce
qui signifie que nous pouvons utiliser des méthodes comme
\texttt{pop\_card} et \texttt{add\_card} pour les mains ainsi que pour
les jeux. Si nous fournissons une méthode \texttt{\_\_init\_\_} dans la
classe \texttt{Hand}, elle remplace celle de la classe \texttt{Deck} :

\begin{Shaded}
\begin{Highlighting}[]
\CommentTok{\# à l\textquotesingle{}intérieur de la classe Hand :}


\KeywordTok{def} \FunctionTok{\_\_init\_\_}\NormalTok{(}\VariableTok{self}\NormalTok{, label}\OperatorTok{=}\StringTok{""}\NormalTok{):}
    \VariableTok{self}\NormalTok{.cards }\OperatorTok{=}\NormalTok{ []}
    \VariableTok{self}\NormalTok{.label }\OperatorTok{=}\NormalTok{ label}
\end{Highlighting}
\end{Shaded}

\hypertarget{diagrammes-de-classes}{%
\subsection{Diagrammes de classes}\label{diagrammes-de-classes}}

Un \textbf{diagramme de classes} est une représentation plus abstraite
de la structure d\textquotesingle un programme qui montre les classes et
leurs relations :

\begin{itemize}
\tightlist
\item
  \textbf{Relation HAS-A (a un)} : les objets d\textquotesingle une
  classe contiennent des références à des objets d\textquotesingle une
  autre classe (par exemple, chaque \texttt{Deck} contient des
  références à de nombreuses \texttt{Cards}).
\item
  \textbf{Relation IS-A (est un)} : une classe hérite
  d\textquotesingle une autre (par exemple, un \texttt{Hand} est une
  sorte de \texttt{Deck}).
\end{itemize}

La multiplicité (comme une étoile \texttt{*} près de la tête de flèche)
indique le nombre d\textquotesingle instances concernées dans une
relation.

\hypertarget{duxe9bogage}{%
\subsection{Débogage}\label{duxe9bogage}}

L\textquotesingle héritage peut représenter un défi pour le débogage car
lorsqu\textquotesingle une méthode est invoquée sur un objet, il
n\textquotesingle est parfois pas évident de savoir quelle
implémentation sera exécutée. La fonction \texttt{find\_defining\_class}
permet de trouver la classe qui fournit la définition
d\textquotesingle une méthode en inspectant l\textquotesingle ordre de
résolution des méthodes (MRO).

\hypertarget{encapsulation-des-donnuxe9es}{%
\subsection{Encapsulation des
données}\label{encapsulation-des-donnuxe9es}}

De la même manière que nous avons découvert des interfaces de fonctions
par l\textquotesingle encapsulation et la généralisation, nous pouvons
découvrir des interfaces de classes par
l\textquotesingle{}\textbf{encapsulation des données}. Le processus de
refactorisation consiste à transformer un ensemble de variables globales
et de fonctions associées en attributs et méthodes d\textquotesingle une
nouvelle classe.

\hypertarget{glossaire}{%
\subsection{Glossaire}\label{glossaire}}

encoder encode Représenter un ensemble de valeurs en utilisant un autre
ensemble de valeurs en construisant une correspondance entre eux.

attribut de classe class attribute Un attribut associé à un objet
classe, défini à l\textquotesingle intérieur d\textquotesingle une
définition de classe mais en dehors de toute méthode.

attribut d\textquotesingle instance instance attribute Un attribut
associé à une instance spécifique d\textquotesingle une classe.

placage veneer Une méthode ou une fonction qui fournit une interface
différente pour une autre fonction sans effectuer de calculs complexes.

héritage inheritance La capacité de définir une nouvelle classe qui est
une version modifiée d\textquotesingle une classe définie précédemment.

classe parente parent class La classe à partir de laquelle une classe
enfant hérite.

classe enfant child class Une nouvelle classe créée en héritant
d\textquotesingle une classe existante ; également appelée « sous-classe
».

relation IS-A IS-A relationship La relation entre une classe enfant et
sa classe parente.

relation HAS-A HAS-A relationship La relation entre deux classes où les
instances d\textquotesingle une classe contiennent des références à des
instances de l\textquotesingle autre.

diagramme de classes class diagram Un diagramme qui montre les classes
d\textquotesingle un programme et les relations entre elles.

multiplicité multiplicity Une notation dans un diagramme de classes qui
indique, pour une relation HAS-A, le nombre de références vers des
instances d\textquotesingle une autre classe.
